\begin{center}
	\zihao{4}\heiti
	%
	\LARGE{ 中国科学院大学\\ $\boldsymbol{2017}$年招收攻读硕士学位研究生入学统一考试试卷\\
	科目名称：电动力学}

\end{center}
\noindent
考生须知：\\
1.本试卷满分为150分,全部考试时间总计180分钟。\\
2.所有答案必须写在答题纸上,写在试题纸上或草稿纸上一律无效。

\vspace*{2ex}
\hrule
\vspace*{4ex}

\fancypagestyle{mypagestyle}{%
	\fancyhf{}% Clear header/footer
	\fancyfoot[OR]{第~\thepage~页，共2页}%
	\fancyfoot[EL]{第~\thepage~页，共2页}%
	\fancyfoot[OL]{科目名称：电动力学 }%
	\fancyfoot[ER]{科目名称：电动力学 }%
	\renewcommand{\headrulewidth}{0pt}% Header rule of .4pt
	\renewcommand{\footrulewidth}{0.5pt}% Header rule of .4pt
}
\fancypagestyle{plain}{%
	\fancyhf{}% Clear header/footer
	\fancyfoot[OR]{第~\thepage~页，共2页}%
	\fancyfoot[EL]{第~\thepage~页，共2页}%
	\fancyfoot[OL]{科目名称：电动力学 }%
	\fancyfoot[ER]{科目名称：电动力学 }%
}

\setcounter{page}{1}%设置页面计数器
\pagestyle{mypagestyle}

\onehalfspacing
% 排版结束，下面是题目

\vspace{0.4cm}

\noindent
一、简答题（共40分）
\vspace{0.2cm}

\noindent
1.\hspace{0.8em}写出无电荷电流分布真空中麦克斯韦方程组，写出电磁场能量密度与能流密度表达式
(8分)\vspace{0.2cm}

\noindent
2.\hspace{0.8em}电偶极子$\Vec{p}$处在原点，写出空间位置$\Vec{r}$处的电势表达式，由此算出其产生的电场强度。(10分)\vspace{0.2cm}

\noindent
3.\hspace{0.8em}写出电磁场矢势和标势的规范变换，用矢势和标势表示电磁场$\Vec{E}, \Vec{B}$，并简要证明$\Vec{E}, \Vec{B}$在此规范变化下不变。(10分)\vspace{0.2cm}

\noindent
4.\hspace{0.8em}写出狭义相对论中粒子的能量-动量关系式。若静止质量为$m_0$的相对论性粒子具有动能$2m_0 c^2$，则此粒子的速度是多少？(12分)\vspace{1cm}

\noindent
二、\hspace{0.2em}单项选择题（共30分，每小题5分）
\vspace{0.2cm}

\noindent
1.\hspace{0.8em}在半径为R的中性导体球外放置一个点电荷，该电荷到球心的距离为2R。现通过一根导线使导体球接地，与接地前相比，此时点电荷所受力：\vspace{0.2cm}

A.\hspace{0.8em}不变  \hspace{16em} B.\hspace{0.8em}方向改变，大小增加到约四倍

C.\hspace{0.8em}方向不变，大小增加到约两倍 \hspace{5em}  D.\hspace{0.8em}方向不变，大小减小到约一半\vspace{0.2cm}

\noindent
2.\hspace{0.8em}某材料能有效屏蔽电磁波，主要原因是\vspace{0.2cm}

A.\hspace{0.8em}介电系数大 \hspace{1em} 
B.\hspace{0.8em}磁导率大 \hspace{1em} 
C.\hspace{0.8em}介电系数与磁导率之积大 \hspace{1em} 
D.\hspace{0.8em}电导率大\vspace{0.2cm}

\noindent
3.\hspace{0.8em}频率为的3kHz平面电磁波垂直入射到电阻率约$10^{-6} \Omega \cdot \mathrm{m}$，磁导率约$10^{-6} \mathrm{H} \cdot \mathrm{m^{-1}} $的某金属表面，其穿透深度约为（准确到数量级）\vspace{0.2cm}

A.\hspace{0.8em}10cm \hspace{1em} 
B.\hspace{0.8em}1cm \hspace{1em} 
C.\hspace{0.8em}0.1cm \hspace{1em} 
D.\hspace{0.8em}0.01cm\vspace{0.2cm}

\noindent
4.\hspace{0.8em}关于电磁波的全反射现象，下面说法正确的是\vspace{0.2cm}

A.\hspace{0.8em}可以发生在任意两种物质的界面处

B.\hspace{0.8em}临界角取决于两种介质的折射率以及电磁波的频率

C.\hspace{0.8em}折射波在界面法向上的平均能流密度为零

D.\hspace{0.8em}反射波与入射波在任意时刻都等幅，反相\vspace{0.2cm}

\noindent
5.\hspace{0.8em}一个圆环一半带均匀正电，另一半带等量的均匀负电。现将其沿着正负电荷交界的中轴线旋转，转动角速度随时间呈正弦振荡，则该体系\vspace{0.2cm}

A.\hspace{0.8em}无辐射 \hspace{16em} 
B.\hspace{0.8em}无电偶极辐射，有磁偶极辐射

C.\hspace{0.8em}有电偶极辐射和电四级辐射 \hspace{7em} 
D.\hspace{0.8em}只有电偶极辐射\vspace{0.2cm}

\noindent
6.\hspace{0.8em}电磁波产生辐射压力的根本原因是因为电磁场具有\vspace{0.2cm}

A.\hspace{0.8em}能量密度 \hspace{1em} B.\hspace{0.8em}动量密度 \hspace{1em} C.\hspace{0.8em}角动量密度 \hspace{1em}  D.\hspace{0.8em}能流密度

\newpage
\noindent
三、\hspace{0.2em}（共30分）有两个同心导体球壳，内球壳半径为$a$，外球壳半径为$b$，分别带电荷$+Q$，$-Q$。两球壳之间的缝隙有一半填充着介电常数为$\varepsilon$的电介质，另一半为真空，球壳系统内外均为真空\vspace{0.2cm}

\noindent
1、写出球壳之间空隙中介质表面处的电场强度和电位移矢量满足的边界条件

\noindent
2、计算空隙中各处的电场强度

\noindent
3、计算两球壳之间的电势差

\noindent
4、计算内球壳面上的球面电荷密度

\noindent
5、计算内球壳处电介质表面上的感生的极化电荷密度\vspace{2cm}

\vspace{-4.5cm}
\begin{figure}[h!]
    \flushright
    \includegraphics[height=4cm]{test/2017/pic_1.png}
\end{figure}


\noindent
四、\hspace{0.2em}（共20分）如图所示，两块无限大的平行薄导体板$A$和$B$上通有垂直纸面向外的均匀分布的电流，电流密度分别为$i_1$和$i_2$，$A$和$B$板之间的区域\uppercase\expandafter{\romannumeral2}充满磁导率为$\mu$的均匀磁介质，区域\uppercase\expandafter{\romannumeral1} 和 \uppercase\expandafter{\romannumeral3}为真空。\vspace{0.2cm}

\noindent
1、求\uppercase\expandafter{\romannumeral1}、\uppercase\expandafter{\romannumeral2}、\uppercase\expandafter{\romannumeral3}三个区域内磁感应强度$B_1$，$B_2$，$B_3$的大小并指明其方向；

\noindent
2、求介质两个面上的磁化电流密度$i_A$和$i_B$并指明磁化电流的方向。\vspace{3cm}

\vspace{-3.5cm}
\begin{figure}[h!]
    \flushright
    \includegraphics[height=3cm]{test/2017/pic_2.png}
\end{figure}

\noindent
五、\hspace{0.2em}（共30分）如图，一束线偏振光垂直入射到平行四边形棱镜的侧边，经两底边的两次反射后，垂直从另一侧边出射。已知棱镜的折射率为n(n>1)，空气折射率为1，棱镜底角为$\phi$。对于第一次底边反射，将该边入射光平行，垂直于入射面的电场分量分别记为$E_{\parallel}$、$E_{\perp}$，相应的磁场分量分别记为$H_{\parallel}$、$H_{\perp}$。将反射光电场磁场的平行，垂直分量分别记为$E_{\parallel}^{'}$、$E_{\perp}^{'}$、$H_{\parallel}^{'}$、$H_{\perp}^{'}$。将折射光电场磁场的平行，垂直分量分别记为$E_{\parallel}^{''}$、$E_{\perp}^{''}$、$H_{\parallel}^{''}$、$H_{\perp}^{''}$。将入射角和折射角记为$\theta$和$\theta^{''}$。
\vspace{0.2cm}

\noindent
1、写出底边上电场磁场满足的边值关系；\vspace{0.1cm}

\noindent
2、当底边反射为全反射时，求棱镜底角满足的条件；\vspace{0.1cm}

\noindent
3、若第一次底边反射为全反射，记$\displaystyle{ \frac{E_{\perp}^{'}}{E_{\perp}} = e^{i \varphi_{\perp}} }$，$\displaystyle{ \frac{E_{\parallel}^{'}}{E_{\parallel}} = e^{i \varphi_{\parallel}} }$，证明$E_{\parallel}^{'}$与$E_{\perp}^{'}$的相位差$\varphi \equiv \varphi_{\parallel} - \varphi_{\perp}$满足$\tan \displaystyle{\frac{\varphi}{2}} = - \displaystyle{  \frac{\cos \sqrt{\sin ^2 \phi - (1/n)^2} }{ \sin ^2 \phi }  } $；\vspace{0.2cm}

\noindent
4、在什么条件下经两次全反射后的出射光为圆偏振光？

\vspace{-1cm}
\begin{figure}[h!]
    \flushright
    \includegraphics[height=3.5cm]{test/2017/pic_3.png}
\end{figure}